% Generated draft: requires mathematical and source review.
\documentclass[11pt]{article}
\usepackage[margin=25mm]{geometry}
\usepackage{fontspec}
\setmainfont{Latin Modern Roman}
\usepackage{amsmath,amssymb,mathtools,unicode-math}
\setmathfont{Latin Modern Math}
\usepackage{enumitem,needspace,xurl,hyperref}
\setlength{\parindent}{0pt}
\setlength{\parskip}{5pt}
\setlength{\emergencystretch}{3em}
\newcommand{\N}{\mathbb{N}}
\newcommand{\Z}{\mathbb{Z}}
\newcommand{\Q}{\mathbb{Q}}
\newcommand{\R}{\mathbb{R}}
\newcommand{\C}{\mathbb{C}}
\newcommand{\F}{\mathbb{F}}
\newcommand{\A}{\mathbb{A}}
\newcommand{\PP}{\mathbb{P}}
\newcommand{\E}{\mathbb{E}}
\newcommand{\eps}{\varepsilon}
\newcommand{\dd}{\,\mathrm{d}}
\begin{document}
\section*{UnsolvedMath ID: 30003120}
\textbf{Source problem code:} catalogue entry
\subsection{Definitions and mathematical statement}
Let \(K(X)\) denote the field of rational functions in an indeterminate \(X\) over a field \(K\). The supplied record does not specify \(K\), its characteristic, or any additional hypotheses on it. For a nonzero rational function \(r=A/B\), with coprime \(A,B\in K[X]\), its degree is \(\deg(r)=\max\{\deg A,\deg B\}\); constant rational maps have degree zero. Composition \(u\circ v\) means substitution \(u(v(X))\), whenever this is defined.

The source asks the following uniform factorization question. Given \(u,f\in K(X)\), suppose that
\[
\exists n\in\mathbb Z_{>0}\;\exists v\in K(X)
\qquad u\circ v=f^n.
\]
Must there exist \(w\in K(X)\) such that
\[
 u\circ w=f^{N},
\]
where the positive integer \(N\) depends only on \(\deg(u)\)? Under a literal uniform-exponent interpretation, this asks for an exponent \(N(d)\in\mathbb Z_{>0}\), chosen uniformly for all admissible \(u,f\) with \(\deg(u)=d\), such that
\[
\left(\exists n\in\mathbb Z_{>0}\;\exists v\in K(X):u\circ v=f^n\right)
\Longrightarrow
\left(\exists w\in K(X):u\circ w=f^{N(d)}\right).
\]
This is a conditional formulation of the source question, not a fully specified assertion over all fields.

\textbf{Unresolved qualifications.} The notation \(f^n\) is not defined in the supplied record. It could mean the ordinary power \((f(X))^n\), or the compositional iterate \(f^{\circ n}\), defined by \(f^{\circ1}=f\) and \(f^{\circ(n+1)}=f\circ f^{\circ n}\) when these compositions are defined. These readings give different questions, and neither is selected here. The record also does not state whether \(u,f,v,w\) must be nonconstant or satisfy degree restrictions. Constant cases require clarification, including whether all displayed substitutions are defined. Furthermore, a uniform exponent \(N(d)\) and an exponent allowed to vary subject to \(1\leq N\leq B(d)\) are different formulations; the question uses the former wording, whereas the supplied literature assessment refers to a degree-dependent bound.

The broader question asks when a reduction to irreducible fibered products of genus at most one produces multiple towers, including more complicated examples. For nonconstant rational maps \(a,b:\mathbb P^1_K\to\mathbb P^1_K\), their fibered product is
\[
 C_{a,b}=\mathbb P^1_K\times_{\mathbb P^1_K}\mathbb P^1_K,
\]
with the two structure maps given by \(a\) and \(b\); its points satisfy \(a(x)=b(y)\). However, the record does not identify the maps used in the reduction, define a tower or distinct towers, specify the transition maps, or state the desired classification. It also does not specify the base field for irreducibility or the precise genus convention. Thus the broader question cannot be sharpened into a uniquely determined mathematical statement from this record alone.

\textbf{Supplied status qualification.} The registry labels the overall problem open. Its dated curation attributes a later proof of the special factorization question, with a degree-dependent bound, to Hyde--Zieve, while describing the broader classification as open. These are supplied assessments, not independently verified mathematical or literature conclusions.
\subsection{Short English statement}
The special question asks whether a factorization \(u\circ v=f^n\) forces a similar factorization at an exponent controlled only by \(\deg(u)\). The meaning of powers, the field hypotheses, constant cases, and the precise uniformity remain unspecified. A broader question concerns multiple towers arising from genus-at-most-one fibered products, but the towers and reduction are not defined. Supplied dated curation reports the special question proved by Hyde--Zieve and the broader question open; this statement-only draft requires human mathematical review.
\subsection{Sources}
\begin{itemize}
\item Local UnsolvedMath registry, supplied complete source record: permanent ID \(30003120\), exact source problem code \texttt{catalogue entry}, title \emph{Uniform Composition Powers of Rational Functions}, collection \emph{Oberwolfach Reports Open Problems}. The record gives proposed year 2016 and supplies no named proposer. Its labels \texttt{open} and \texttt{corrected\_verified} are registry metadata, not independent verification here.
\item Supplied report citation: \emph{Diophantische Approximationen (2017)}, Oberwolfach Reports, pp.\ 1099--1187, DOI \texttt{10.4171/owr/2016/21}. \url{https://doi.org/10.4171/owr/2016/21}. The record describes these reports as documenting mathematical workshops and their open-problem discussions.
\item Supplied literature source, labelled ``OWR report'': \url{https://ems.press/content/serial-article-files/46625?nt=1}.
\item Supplied literature source, labelled ``Hyde thesis'': \url{https://pages.vassar.edu/thyde/files/2025/11/Hyde-Thesis-ADML.pdf}. The record attributes the later special-factorization result to Hyde--Zieve but gives no theorem number or separate bibliographic citation for that result.
\item Literature assessment date: \texttt{2026-08-21}. This is supplied dated curation, described in the record as literature triage derived from the supplied OWR curation, and has not been independently verified. It reports a degree-dependent bound for the special factorization question and an open broader classification of multiple towers and more complicated examples. The record explicitly calls for independent verification of time-sensitive status and cited claims.
\end{itemize}
\end{document}
